\documentclass[11pt,letterpaper]{article}
\usepackage[margin=1in]{geometry}
\usepackage{amsmath,amssymb,amsthm}
\usepackage{hyperref}
\usepackage{booktabs}
\usepackage{listings}
\usepackage{xcolor}
\usepackage{enumitem}
\usepackage{microtype}
\usepackage{graphicx}
\usepackage{longtable}
\usepackage{array}
\usepackage{tabularx}

\hypersetup{colorlinks=true,linkcolor=blue,citecolor=blue,urlcolor=blue}

\lstset{
  language=Fortran,
  basicstyle=\small\ttfamily,
  keywordstyle=\bfseries\color{blue!70!black},
  commentstyle=\itshape\color{gray},
  stringstyle=\color{orange!80!black},
  numbers=left,
  numberstyle=\tiny\color{gray},
  frame=single,
  breaklines=true,
  captionpos=b
}

\title{\textbf{Liquid-Metal MHD in MUG/OFT:\\
Capabilities, Comparison, Improvement Roadmap,\\
and Multi-Physics Coupling Game Plan}\\
\large OpenFUSIONToolkit --- Commit \texttt{5faae9a} and beyond}

\author{Prepared for: Sophia Guizzo and OFT Maintainers\\
Compiler: Claude Code (repository-grounded analysis)}

\date{Repository state: \texttt{main}, HEAD \texttt{5faae9ad698c2014afbe35b4ab837a6d1482b6b1}\\
Analysis date: 2026-05-08}

\begin{document}
\maketitle
\tableofcontents
\newpage

%==============================================================================
\section{Abstract}
%==============================================================================

The OpenFUSIONToolkit (OFT) MHD solver MUG has recently been extended with a
2D/axisymmetric module (\texttt{xmhd\_2d.F90}, commit \texttt{5faae9a}).  This
module solves the compressible, resistive, viscous MHD equations in two
spatial dimensions (either Cartesian or cylindrical $R$--$Z$) for seven coupled
fields: number density $n$, velocity components $(v_x, v_y, v_z)$, temperature
$T$, in-plane magnetic flux $\psi$, and out-of-plane field $B_y$.  The
implementation is complete for basic plasma physics and is validated against
Alfv\'{e}n and sound-wave benchmarks.

According to a recent APS DPP poster (Guizzo, Columbia University), an
axisymmetric MUG implementation has been developed for coupled plasma and
liquid-metal blanket dynamics, with TokaMaker providing the evolving plasma
equilibrium.  This poster capability maps partially onto the current
\texttt{main} branch: TokaMaker has a one-way equilibrium export
(\texttt{save\_mug()}) and a point-wise field interpolator, while
\texttt{xmhd\_2d.F90} has full cylindrical $R$--$Z$ support.  However,
liquid-metal-specific physics (wall drag, Hartmann friction, conducting-wall
BCs, low-Rm formulation), the plasma--blanket feedback path, and a MUG Python
interface are all absent from \texttt{main}.

This document (1) surveys what MUG can already do, (2) maps the APS poster
capabilities onto the repository, (3) compares MUG to peer liquid-metal MHD
codes, (4) presents a physics improvement roadmap ranked by impact and effort,
(5) analyses the existing and missing multi-physics coupling infrastructure, and
(6) proposes a sequenced game plan of pull requests from first drag term to
full plasma--blanket coupling.

%==============================================================================
\section{Executive Summary}
%==============================================================================

\begin{enumerate}[label=\textbf{ES-\arabic*.}]
  \item \texttt{xmhd\_2d.F90} (commit \texttt{5faae9a}) provides a
    production-quality 2D MHD solver with uniform scalar viscosity $\nu$,
    resistivity $\eta$, thermal diffusivity $\chi$, and particle diffusivity
    $D$.  Cylindrical $R$--$Z$ coordinates are enabled via \texttt{cyl\_flag}
    throughout all hot paths.

  \item \textbf{No wall-drag, Hartmann-friction, Shercliff-layer, traction BC,
    or liquid-metal-specific terms exist anywhere in the repository.}  The APS
    poster feature is partially visible (cylindrical flag, TokaMaker equilibrium
    export), but the wall-physics code, material-specific closures, and
    plasma--blanket feedback loop are absent from \texttt{main}.

  \item Compared to peer codes, MUG's 2D module is solving harder equations
    than necessary for liquid metals: the full compressible induction equation
    with Alfv\'{e}nic timestep constraint, when Rm $\ll 1$ and Ma $\ll 1$ for
    all engineering liquid-metal conditions.  The most impactful single upgrade
    is a \textbf{low-Rm electric potential solver mode} that replaces the
    induction equation with a Poisson solve and allows timesteps 10$^2$--10$^3\times$
    larger.

  \item \textbf{Seven physics improvements} can be implemented without
    architectural changes, ordered by effort-to-impact:
    (1)~gravity/buoyancy, (2)~ohmic+viscous heating, (3)~density/temperature
    floors, (4)~SM82 Hartmann+Shercliff friction from $\mathrm{Ha}$, (5)~conducting-wall
    Robin BC on $\psi$, (6)~low-Rm electric potential mode, (7)~port of 3D
    anisotropic viscosity and regional resistivity.

  \item The multi-physics coupling infrastructure is \textbf{partially present}:
    TokaMaker has a one-way HDF5 equilibrium export to MUG and a live field
    interpolator; ThinCurr has cross-model electromagnetic coupling; the
    TokaMaker--TORAX temporal coupling loop in \texttt{pulse\_design.py} is the
    architectural template.  The missing pieces are: a MUG Python API, a live
    $\mathbf{B}_0$ injection path from TokaMaker's interpolator into MUG, and a
    $\mathbf{J}_{\mathrm{blanket}} \to$ TokaMaker feedback interface.

  \item The \textbf{recommended game plan is eight sequenced PRs} from the
    minimal Hartmann drag term (PR~1) through full bidirectional plasma--blanket
    coupling (PR~8), with each step independently useful and testable.
\end{enumerate}

%==============================================================================
\section{Repository and Commit State (Code-Confirmed Facts)}
%==============================================================================

\subsection{Git Provenance}

\begin{center}
\begin{tabular}{ll}
\toprule
\textbf{Item} & \textbf{Value} \\
\midrule
Branch & \texttt{main} \\
HEAD & \texttt{5faae9ad698c2014afbe35b4ab837a6d1482b6b1} \\
Parent & \texttt{5e6ea6283ba0b88690df59859b4cafce63d04088} \\
Commit type & \texttt{commit} (verified via \texttt{git cat-file -t 5faae9a}) \\
Working tree & Clean (only \texttt{CLAUDE.MD} untracked) \\
\bottomrule
\end{tabular}
\end{center}

\subsection{Commit 5faae9a: Files Added and Modified}

Commit \texttt{5faae9a} (``MUG: Add 2D implementation \#216'') was authored by
samuel-frei and co-authored by Sophia Guizzo (\texttt{@sguizzo}), Shreyas
Seethalla (\texttt{@KasaiYuki}), and Chris Hansen (\texttt{@hansec}).  The
diff introduced 3{,}439 net insertions across 21 files.

\textbf{New files added:}
\begin{itemize}
  \item \texttt{src/physics/xmhd\_2d.F90} --- 1{,}812 lines, the 2D MHD module
  \item \texttt{src/examples/MUG/harris\_sheet/\{CMakeLists.txt, harris\_sheet.F90, oft.in, oft\_in.xml\}}
  \item \texttt{src/tests/physics/\{test\_alfven\_2d.F90, test\_alfven\_2d.py, test\_sound\_2d.F90, test\_sound\_2d.py\}}
\end{itemize}

\textbf{Modified files (significant):}
\begin{itemize}
  \item \texttt{src/base/oft\_io.F90} (+55) --- I/O extensions
  \item \texttt{src/fem/blag\_operators.F90} (+44) --- new Lagrange operators
  \item \texttt{src/fem/fem\_utils.F90} (+28) --- FEM utilities
  \item \texttt{src/physics/diagnostic.F90} (+29) --- diagnostic output
  \item \texttt{src/physics/oft\_scalar\_inits.F90} (+68) --- scalar initialisation
  \item \texttt{src/lin\_alg/\{lin\_alg\_base.F90, native\_la.F90\}} --- linear algebra
  \item \texttt{src/grid/mesh\_cube.F90} (+15) --- structured mesh support
\end{itemize}

%==============================================================================
\section{APS DPP Poster Context (Poster-Sourced Facts)}
%==============================================================================

The APS DPP poster titled ``Flexible, Axisymmetric MHD Solver for Coupled
Plasma and Liquid Metal Dynamics'' (presenter: Sophia Guizzo, Columbia
University) reports the following capabilities:

\begin{enumerate}
  \item An \textbf{axisymmetric} implementation of MUG within OFT solving the
    \textbf{visco-resistive compressible MHD equations} with high geometric
    fidelity.
  \item Coupling of axisymmetric MUG with \textbf{TokaMaker}, the 2D
    time-dependent plasma equilibrium solver.
  \item The coupling computes \textbf{current distribution and flow profile in
    the blanket} due to time-evolving magnetic fields.
  \item Blanket response \textbf{feeds back on concurrent plasma evolution}.
  \item Application to \textbf{disruption impact on blanket fluid} and whether
    blanket response alters plasma dynamics.
\end{enumerate}

\textbf{Mapping to repository (code-confirmed):}

\begin{center}
\begin{longtable}{p{0.45\textwidth}p{0.5\textwidth}}
\toprule
\textbf{Poster Claim} & \textbf{Repository Status} \\
\midrule
Axisymmetric MUG & \textbf{Confirmed partially.}  \texttt{cyl\_flag} in
  \texttt{oft\_xmhd\_2d\_sim} (line~71) enables cylindrical $R$--$Z$
  coordinates throughout \texttt{nlfun\_apply} and \texttt{build\_approx\_jacobian}. \\
\addlinespace
Visco-resistive compressible MHD & \textbf{Confirmed.} $\nu$, $\eta$, $\chi$,
  $D$, $\gamma$ are all implemented. \\
\addlinespace
TokaMaker coupling (one-way) & \textbf{Confirmed.} \texttt{save\_mug()} in
  \texttt{TokaMaker/\_core.py} (line~1905) exports equilibrium (mesh, $B_R$,
  $B_T$, $B_Z$, $P$, $\psi$) as HDF5 via \texttt{gs\_save\_mug} in
  \texttt{grad\_shaf.F90} (lines 4851--4939). \\
\addlinespace
TokaMaker live field access & \textbf{Confirmed.} \texttt{TokaMaker\_field\_interpolator}
  (\texttt{\_core.py} line~2606) evaluates $\mathbf{B}$, $\psi$, $P$, $\nabla\mathbf{B}$
  at arbitrary $(R,Z)$ points. \\
\addlinespace
Blanket material properties & \textbf{Not found.} No liquid-metal density,
  conductivity, or viscosity parameters specific to PbLi or liquid lithium. \\
\addlinespace
Current distribution in blanket & \textbf{Not found.} No region-specific
  conductivity or blanket geometry in \texttt{xmhd\_2d.F90}. \\
\addlinespace
Plasma--blanket feedback loop & \textbf{Not found.} No coupled simulation
  driver in \texttt{main}; feedback path MUG $\to$ TokaMaker is absent. \\
\addlinespace
Disruption transients & \textbf{Not found.} No disruption-specific examples or
  input files for liquid-metal channels. \\
\bottomrule
\end{longtable}
\end{center}

\textbf{Conclusion:} The axisymmetric visco-resistive MHD infrastructure is
present; the one-way TokaMaker equilibrium export works.  Liquid-metal
material closures, wall-drag physics, and the plasma--blanket feedback loop
are absent from \texttt{main} or unreleased.

%==============================================================================
\section{Current MUG 2D/Axisymmetric Architecture (Code-Confirmed Facts)}
%==============================================================================

\subsection{Module Overview}

File: \texttt{src/physics/xmhd\_2d.F90} (1{,}812 lines)\\
Module name: \texttt{xmhd\_2d}\\
FEM basis: Lagrange scalar elements (\texttt{oft\_blagrange}) on 2D boundary meshes

\subsection{Primary Derived Type: \texttt{oft\_xmhd\_2d\_sim}}

Lines 69--124.  Key fields:

\begin{center}
\begin{tabular}{llll}
\toprule
\textbf{Field} & \textbf{Default} & \textbf{Units} & \textbf{Line} \\
\midrule
\texttt{nu} & $-1$ & m$^2$/s (kinematic viscosity) & 83 \\
\texttt{eta} & --- & $\mu_0 \cdot \Omega\cdot\text{m}$ (resistivity) & 92 \\
\texttt{chi} & $-1$ & m$^2$/s (thermal diffusivity) & 82 \\
\texttt{D\_diff} & $-1$ & m$^2$/s (particle diffusivity) & 85 \\
\texttt{gamma} & $2/3$ & --- (adiabatic index) & 84 \\
\texttt{den\_scale} & $10^{19}$ & m$^{-3}$ (density normalisation) & 87 \\
\texttt{m\_i} & $m_p$ & kg (ion mass) & 88 \\
\texttt{k\_boltz} & $e$ (eV mode) & J/K & 86 \\
\texttt{B\_0(3)} & $[0,0,0]$ & T (background field) & 91 \\
\texttt{cyl\_flag} & \texttt{.FALSE.} & --- & 71 \\
\texttt{linear} & \texttt{.FALSE.} & --- & 72 \\
\texttt{timestep\_cn} & \texttt{.FALSE.} & --- & 73 \\
\texttt{dt} & $-1$ & s & 79 \\
\texttt{nsteps} & $-1$ & --- & 74 \\
\texttt{nl\_tol} & $10^{-5}$ & --- & 90 \\
\texttt{lin\_tol} & $10^{-8}$ & --- & 89 \\
\bottomrule
\end{tabular}
\end{center}

\subsection{Evolved Fields (7 total)}

The solution vector $\mathbf{u}$ contains seven scalar fields at every Lagrange
node:

\begin{center}
\begin{tabular}{cll}
\toprule
\textbf{Index} & \textbf{Field tag} & \textbf{Physical meaning} \\
\midrule
1 & \texttt{n} & number density \\
2 & \texttt{velx} & $v_x$ (or $v_R$ in cylindrical) \\
3 & \texttt{vely} & $v_y$ (or $v_Z$ in cylindrical) \\
4 & \texttt{velz} & $v_z$ (toroidal/out-of-plane) \\
5 & \texttt{T} & temperature \\
6 & \texttt{psi} & in-plane magnetic flux function \\
7 & \texttt{by} & out-of-plane magnetic field \\
\bottomrule
\end{tabular}
\end{center}

The velocity triplet $(v_x, v_y, v_z)$ on a 2D mesh gives a \textbf{2.5D}
formulation: all fields depend on only two spatial coordinates, but
three vector components are retained.

\subsection{Governing Equations}

The module solves (in Cartesian notation; cylindrical adds metric factors at
\texttt{cyl\_flag} branches):

\begin{align}
\frac{\partial n}{\partial t} + \nabla\cdot(n\mathbf{v}) &= D\nabla^2 n
\label{eq:continuity} \\
m_i n\left(\frac{\partial \mathbf{v}}{\partial t}
  + \mathbf{v}\cdot\nabla\mathbf{v}\right)
  &= \mathbf{J}\times\mathbf{B}
    - \nabla(2nk_BT)
    + \nu\nabla^2\mathbf{v}
    + \nu\frac{\nabla n}{n}\cdot\nabla\mathbf{v}
\label{eq:momentum} \\
\frac{n}{\gamma-1}\left(\frac{\partial T}{\partial t}
  + \mathbf{v}\cdot\nabla T\right)
  &= -nT\nabla\cdot\mathbf{v}
    + \chi\nabla^2 T
    - \chi\frac{\nabla n}{n}\cdot\nabla T
\label{eq:energy} \\
\frac{\partial\psi}{\partial t} + \mathbf{v}\cdot\nabla\psi
  &= \eta\nabla^2\psi - (\mathbf{B}_0\times\mathbf{v})_y
\label{eq:psi} \\
\frac{\partial B_y}{\partial t}
  &= \eta\nabla^2 B_y + [\nabla\times(\nabla\psi\times\mathbf{v})]_y
\label{eq:by}
\end{align}

\subsection{3D XMHD: Physics Present in 3D But Missing in 2D}

The 3D \texttt{xmhd.F90} has all of the following, none of which have been
ported to \texttt{xmhd\_2d.F90}:

\begin{itemize}
  \item \texttt{nu\_par}, \texttt{nu\_perp} --- anisotropic (field-aligned) viscosity
  \item \texttt{eta\_reg(:)} --- per-region resistivity scale factors
  \item \texttt{eta\_temp} --- Spitzer-profile temperature-dependent resistivity
  \item \texttt{eta\_hyper} --- hyper-resistivity
  \item Hall term: $(1/ne)(\mathbf{J}\times\mathbf{B} - \nabla nkT_e)$
  \item Two-temperature model (\texttt{xmhd\_two\_temp})
  \item Braginskii anisotropic thermal conductivity (\texttt{xmhd\_brag})
  \item Ohmic heating $\eta J^2$ (\texttt{xmhd\_ohmic})
  \item Viscous heating (\texttt{xmhd\_visc\_heat})
  \item Gravity (\texttt{g\_accel})
  \item Density and temperature floors (\texttt{den\_floor}, \texttt{temp\_floor})
\end{itemize}

%==============================================================================
\section{Comparison with Existing Liquid-Metal MHD Codes}
%==============================================================================

\subsection{The Peer Landscape}

\begin{center}
\begin{longtable}{p{0.13\textwidth}p{0.13\textwidth}p{0.12\textwidth}p{0.12\textwidth}p{0.14\textwidth}p{0.27\textwidth}}
\toprule
\textbf{Code} & \textbf{Group} & \textbf{Rm} & \textbf{Compress.} & \textbf{Geometry} & \textbf{Wall physics} \\
\midrule
\textbf{HIMAG} & UCLA (Moreau) & Low-Rm (inductionless, $\varphi$ form.) & Incompressible & 3D unstructured FEM & Conducting-wall BCs, arbitrary $c_w$ \\
\addlinespace
\textbf{HiMag} & UCLA (Smolentsev) & Low-Rm & Incompressible & 3D structured FVM, RANS & Wall functions, Ha $\sim10^4$ \\
\addlinespace
\textbf{PTOLEMY} & Pothérat group & Low-Rm & Incompressible & \textbf{2D quasi-2D (SM82)} & Hartmann + Shercliff friction \\
\addlinespace
\textbf{OpenFOAM-MHD} & Community & Low-Rm ($\varphi$ form.) & Incompressible & 3D unstructured FVM & Basic no-slip \\
\addlinespace
\textbf{SFEMANS} & Guermond--Léorat & Full Rm & Both & Axisymmetric + Fourier & Limited \\
\addlinespace
\textbf{MUG xmhd} & OFT (Hansen) & \textbf{Full Rm} & \textbf{Compressible} & 3D unstructured FEM & No-slip Dirichlet only \\
\addlinespace
\textbf{MUG xmhd\_2d} & OFT (Guizzo) & \textbf{Full Rm} & \textbf{Compressible} & 2D/axisymmetric FEM & No-slip Dirichlet only \\
\bottomrule
\end{longtable}
\end{center}

\subsection{The Two Core Mismatches for Liquid Metals}

\subsubsection{Magnetic Reynolds Number}

For engineering liquid metals, Rm $= \mu_0 \sigma v L$ is very small:

\begin{center}
\begin{tabular}{llll}
\toprule
\textbf{Material} & $\sigma$ (S/m) & $v$ (m/s) & Rm at $L = 0.1$ m \\
\midrule
Liquid lithium (500\,K) & $3.0\times10^6$ & 0.1 & $\approx 4\times10^{-2}$ \\
PbLi (623\,K) & $7.5\times10^5$ & 0.1 & $\approx 9\times10^{-3}$ \\
\bottomrule
\end{tabular}
\end{center}

At Rm $\ll 1$, the magnetic field is barely perturbed by the flow.  Solving
MUG's full induction equations~\eqref{eq:psi}--\eqref{eq:by} is physically
correct but wasteful: the Alfv\'{e}n speed sets the timestep, but the
Alfv\'{e}n crossing time is $\sim 10^3$--$10^5\times$ shorter than the
flow timescale.  This forces far more timesteps than the physics requires.
All peer codes except SFEMANS avoid this by adopting the \textbf{inductionless
(electric-potential) approximation}: the induced field is negligible, so one
solves a scalar Poisson equation for electric potential $\varphi$ instead.

\subsubsection{Mach Number and Compressibility}

Liquid-metal flow speeds are $v \sim 0.01$--$1$\,m/s; the sound speed in
liquid lithium is $c_s \approx 4{,}500$\,m/s, giving Ma $\sim 10^{-4}$.  The
acoustic timestep constraint $\Delta t \lesssim \Delta x / c_s$ forces steps
many orders of magnitude smaller than necessary to resolve the actual flow.
All peer codes use incompressible formulations and avoid this entirely.

\subsection{What MUG Has That Peer Codes Do Not}

\begin{itemize}
  \item \textbf{Full Rm capability:} MUG can handle Alfv\'{e}nic and
    electromagnetic timescales that low-Rm codes cannot.  This is essential
    for the plasma side of the coupled problem (where Rm $\sim 1$--$10$) and
    for regimes where the induced field in the liquid metal is not entirely
    negligible.

  \item \textbf{Unstructured FEM:} Enables complex axisymmetric blanket
    geometries with port openings, curved walls, and field-shaping inserts
    without mapping to a structured grid.

  \item \textbf{Newton--Krylov solver:} Implicit nonlinear time-stepping with
    block-Jacobi preconditioner scales to stiff problems with large variations
    in timescale.

  \item \textbf{Native OFT integration:} Shares mesh, I/O, and basis
    infrastructure with TokaMaker and ThinCurr, enabling future tight coupling
    without format translation.

  \item \textbf{Cylindrical 2.5D:} Out-of-plane field and velocity components
    on a 2D mesh capture toroidal flows and poloidal currents without full 3D,
    which is the natural geometry for a tokamak blanket.
\end{itemize}

\subsection{PTOLEMY as the Closest Analogue}

PTOLEMY (Sommeria--Moreau quasi-2D model) is the best comparison target for
MUG's 2D module.  It is 2D, incompressible, and uses the SM82 model to
represent unresolved Hartmann and Shercliff layers as volume friction.  MUG's
2D module with the improvements in Section~\ref{sec:roadmap} would be:
\begin{itemize}
  \item \textbf{More general:} unstructured mesh vs.\ PTOLEMY's structured grid;
    full 2.5D field evolution vs.\ 2D vorticity stream-function.
  \item \textbf{Equally tractable:} same 2D solve on similar mesh sizes.
  \item \textbf{Plasma-coupled:} directly tied to TokaMaker for the evolving
    background field, which PTOLEMY has no mechanism for.
\end{itemize}

%==============================================================================
\section{Existing Liquid-Metal / Blanket Capabilities (Code-Confirmed Facts)}
%==============================================================================

A systematic keyword search across all \texttt{.F90}, \texttt{.py},
\texttt{.in}, \texttt{.xml}, and \texttt{.md} files returned \textbf{zero
matches} for: \texttt{liquid}, \texttt{lithium}, \texttt{lead},
\texttt{blanket}, \texttt{PbLi}, \texttt{Hartmann}, \texttt{Shercliff},
\texttt{drag}, \texttt{friction}, \texttt{conducting.wall}.

\textbf{Conclusion:} At \texttt{main} HEAD there is no liquid-metal-specific
physics.  MUG's equations are physically correct for a conducting fluid, but no
blanket material parameters, Hartmann or Shercliff friction models,
wall-conductivity treatments, or plasma--blanket coupling loops are implemented.

%==============================================================================
\section{Physics Improvement Roadmap}
\label{sec:roadmap}
%==============================================================================

The following improvements are ordered by the ratio of physics impact to
implementation effort.  None requires architectural changes to MUG's solver
structure; all fit within the existing Newton--Krylov, Lagrange-FEM framework.

\subsection{Group A: Port Physics from 3D \texttt{xmhd.F90} (Low Effort)}

These features are already implemented and validated in the 3D module.
Porting to \texttt{xmhd\_2d.F90} is primarily transcription with cylindrical
metric factors.

\subsubsection{A1. Gravity and Boussinesq Buoyancy ($\sim$20 lines)}

The 3D module has a \texttt{g\_accel} field.  Add to the momentum residual:
\begin{equation}
  \mathbf{S}_g = m_i n \mathbf{g}\left[1 - \beta_T(T - T_{\mathrm{ref}})\right]
\end{equation}
where $\beta_T$ is the thermal expansion coefficient.  This unlocks
magnetoconvection problems: Rayleigh--B\'{e}nard convection in a magnetic
field, natural convection in a heated blanket channel.  Standard validation
benchmarks exist (Burr \& M\"{u}ller 2001, Auth\'{e}i et al.\ 2003).

\subsubsection{A2. Ohmic and Viscous Heating ($\sim$15 lines)}

Add to the temperature residual:
\begin{equation}
  Q_{\mathrm{ohm}} = \eta J^2, \qquad
  Q_{\mathrm{visc}} = \frac{1}{2}(\nabla\mathbf{v})^T : \Pi
\end{equation}
Critical for disruption scenarios where large induced currents deposit ohmic
heat in the blanket.

\subsubsection{A3. Density and Temperature Floors ($\sim$10 lines)}

Add \texttt{den\_floor} and \texttt{temp\_floor} guards to prevent
negative density or temperature near sharp interfaces or shocks.
The 3D module has these; they cost nothing computationally.

\subsubsection{A4. Temperature-Dependent Resistivity ($\sim$20 lines)}

The 3D module implements Spitzer profile $\eta(T) \propto T^{-3/2}$.  For
liquid metals a linear metallic model $\eta(T) = \eta_0(1 + \alpha_\eta(T-T_0))$
is more appropriate.  The infrastructure (user-switchable profile function)
is what matters.

\subsubsection{A5. Anisotropic Viscosity Port ($\sim$30 lines)}

Port \texttt{nu\_par}, \texttt{nu\_perp} from the 3D module.  The viscous
stress tensor becomes:
\begin{equation}
  \Pi = -\left[\nu_\parallel \hat{b}\hat{b}
              + \nu_\perp(I - \hat{b}\hat{b})\right] : W
\end{equation}
Important for magnetoconvection and magnetically aligned flow in blanket channels.

\subsection{Group B: SM82 Quasi-2D Wall Drag (Medium Effort)}
\label{sec:sm82}

The Sommeria--Moreau (1982) quasi-2D model is the standard reduced
representation of unresolved Hartmann and Shercliff boundary layers.  It
replaces thin-layer physics with two volumetric friction terms:

\begin{equation}
  \mathbf{S}_{u}^{(\mathrm{SM82})} =
  -\underbrace{\alpha_H n \mathbf{v}_\perp}_{\text{Hartmann walls}}
  - \underbrace{\alpha_S n \mathbf{v}}_{\text{Shercliff walls}}
\end{equation}

where, in dimensional form:
\begin{align}
  \alpha_H &= \frac{\sigma B^2}{\rho} \cdot \frac{1}{\mathrm{Ha}}
            = \frac{\nu \,\mathrm{Ha}}{L^2}
  \label{eq:alpha_H} \\
  \alpha_S &\approx \alpha_H \cdot \mathrm{Ha}^{-1/2}
            = \frac{\nu \,\mathrm{Ha}^{1/2}}{L^2}
  \label{eq:alpha_S}
\end{align}

and $\mathrm{Ha} = BL\sqrt{\sigma/(\rho\nu)}$.

\textbf{Implementation:} Add three input parameters to \texttt{oft\_xmhd\_2d\_sim}:
\texttt{hartmann\_number} (Ha), \texttt{wall\_spacing} ($L$),
\texttt{B0\_direction} (unit vector parallel to $\mathbf{B}$).  Compute
$\alpha_H$ and $\alpha_S$ at setup.  The residual addition is the same as
Option~A/B in Section~\ref{sec:impl_options} but with physically derived
coefficients and directional selectivity.

\textbf{Comparison:} This is precisely the model used by PTOLEMY.  With this
addition, MUG's 2D module becomes PTOLEMY-comparable for quasi-2D liquid-metal
flows, with the advantage of unstructured mesh and OFT's plasma coupling.

\subsection{Group C: Conducting-Wall Robin BC on $\psi$ (Medium Effort)}

Realistic tokamak blankets have conducting walls.  The wall conductance ratio:
\begin{equation}
  c_w = \frac{\sigma_w t_w}{\sigma_f L}
\end{equation}
modifies the effective Hartmann number and the current closure condition.
The magnetic BC on $\psi$ at a conducting wall changes from Dirichlet
($\psi = 0$, insulating) to Robin:
\begin{equation}
  \psi + c_w \frac{\partial\psi}{\partial n} = 0 \quad \text{on wall}
  \label{eq:robin_psi}
\end{equation}

\textbf{Implementation:} Requires adding a boundary-face quadrature loop and
integrating the traction term in the $\psi$ residual --- the one missing
weak-form infrastructure piece in \texttt{xmhd\_2d.F90}.

\textbf{Validation:} Hunt (1965) analytic solution for duct flow with
conducting walls is the canonical benchmark.  This makes MUG directly
comparable to HIMAG for the conducting-wall case.

\subsection{Group D: Low-Rm Electric Potential Mode (Medium-to-Large Effort, Highest Payoff)}
\label{sec:low_rm}

This is the single most impactful upgrade for liquid-metal physics.
When Rm $\ll 1$:

\begin{enumerate}
  \item The applied field $\mathbf{B}_0$ is prescribed (from TokaMaker or
    analytically).
  \item Solve for electric potential $\varphi$:
    \begin{equation}
      \nabla\cdot(\sigma\nabla\varphi) = \nabla\cdot(\mathbf{v}\times\mathbf{B}_0)
      \label{eq:phi_poisson}
    \end{equation}
  \item Compute current:
    $\mathbf{J} = \sigma(-\nabla\varphi + \mathbf{v}\times\mathbf{B}_0)$
  \item Add Lorentz force $\mathbf{J}\times\mathbf{B}_0$ to momentum.
\end{enumerate}

This \textbf{replaces the stiff induction equations}~\eqref{eq:psi}--\eqref{eq:by}
with one scalar Poisson solve, reducing the field count from 7 to 5 and
removing the Alfv\'{e}nic timestep constraint.  Practical speedup for
liquid-metal conditions: $10^2$--$10^3\times$ larger $\Delta t$.

\textbf{Implementation:} New solver mode \texttt{use\_low\_rm = .TRUE.} in
\texttt{oft\_xmhd\_2d\_sim}.  The $\varphi$ field uses the same Lagrange
FEM basis as the existing fields.  The Poisson operator is simpler than the
existing induction Jacobian.  Boundary condition: Neumann
($\mathbf{J}\cdot\hat{n} = 0$) at insulating walls, or Robin for conducting
walls (same $c_w$ parameter as in Sec.\ C).

\textbf{Coupling advantage:} With this mode, MUG does not evolve its own
magnetic field at all --- it reads $\mathbf{B}_0(t)$ from TokaMaker's field
interpolator at each step.  This is the natural interface for TokaMaker coupling
(Section~\ref{sec:coupling_roadmap}).

\textbf{Comparison:} HIMAG and HiMag both use exactly this formulation.
With Group~D implemented, MUG's 2D module becomes directly comparable to HIMAG
for the inductionless regime.

\subsection{Group E: Regional Material Properties (Larger Effort)}

Port \texttt{eta\_reg(:)} from 3D \texttt{xmhd.F90} and extend to a general
per-element material map for $\nu$, $\eta$, $\rho$, $\sigma$.  This enables
simultaneous simulation of plasma and liquid-metal domains with different
material properties and a physical interface condition.  This is necessary for
the full multi-material blanket geometry but is not required for the initial
coupling studies where MUG is run in the blanket region only.

\subsection{Summary Table}

\begin{center}
\begin{longtable}{p{0.22\textwidth}p{0.08\textwidth}p{0.08\textwidth}p{0.25\textwidth}p{0.27\textwidth}}
\toprule
\textbf{Improvement} & \textbf{Effort} & \textbf{Lines} & \textbf{Physics gain} & \textbf{Comparable to} \\
\midrule
A1. Gravity/buoyancy & Low & $\sim$20 & Magnetoconvection & HIMAG, Burr \& Müller \\
A2. Ohmic+viscous heating & Low & $\sim$15 & Disruption energy balance & 3D xmhd \\
A3. Floors & Low & $\sim$10 & Robustness & 3D xmhd \\
A4. $\eta(T)$ profile & Low & $\sim$20 & Metallic resistivity & 3D xmhd \\
A5. Anisotropic viscosity & Low & $\sim$30 & Magnetically aligned flow & 3D xmhd \\
B. SM82 drag (Ha-based) & Medium & $\sim$60 & Quasi-2D LM flow & PTOLEMY \\
C. Conducting-wall Robin BC & Medium & $\sim$100 & Hunt channel benchmark & HIMAG \\
D. Low-Rm $\varphi$ mode & Medium-large & $\sim$400 & Full LM physics, $10^{2-3}\times \Delta t$ & HIMAG, HiMag \\
E. Regional materials & Large & $\sim$300+ & Multi-material domains & All 3D codes \\
\bottomrule
\end{longtable}
\end{center}

%==============================================================================
\section{Physics Model Hierarchy for Wall Drag}
%==============================================================================

\subsection{Tier 1: Viscous No-Slip Wall Drag}

\textbf{Is it already captured?}  The default BC in \texttt{setup\_bc}
(lines~$\approx$1602--1611) applies zero Dirichlet velocity at all boundary
nodes, enforcing no-slip.  However, the viscous Hartmann layer thickness
$\delta_\nu \sim L/\mathrm{Ha}$ at engineering Ha $\sim 10^3$--$10^4$ is
$10^{-3}$--$10^{-4} L$ --- almost never resolved at blanket scale.

\subsection{Tier 2: Quasi-2D Hartmann Friction}
(Recommended first addition --- see Section~\ref{sec:sm82})

\begin{equation}
  \mathbf{S}_{u}^{(H)} = -\alpha_H\, m_i n\, \mathbf{v}_\perp, \qquad
  \alpha_H = \frac{\sigma B^2}{m_i n}
  \label{eq:hartmann_drag}
\end{equation}

Valid for Ha $\gg 1$, quasi-2D flow.  Modelling Hartmann layers without
resolving them.

\subsection{Tier 3: Shercliff-Layer Side-Wall Drag}

\begin{equation}
  \mathbf{S}_{u}^{(S)} = -\beta_S\, m_i n\, \mathbf{v}, \qquad
  \beta_S \propto \frac{\nu}{L^2}\,\mathrm{Ha}^{1/2}
\end{equation}

Can be combined with Tier~2 into the SM82 model.

\subsection{Tier 4: Conducting-Wall Electromagnetic Coupling}

Wall conductivity $\sigma_w$ modifies the current closure condition
(Robin BC on $\psi$, Eq.~\eqref{eq:robin_psi}) and the effective Ha.
Implemented via \texttt{thin\_wall.F90} in 3D; a 2D analogue is needed.

\subsection{Tier 5: Fully Resolved 3D Hartmann Layers}

Requires mesh spacing $\lesssim \delta_H = L/\mathrm{Ha}$, anisotropic
wall-normal meshing, 3D geometry, and Alfv\'{e}nic timestep.  Tractable only
at low Ha and small domains.  Not a near-term target.

%==============================================================================
\section{Does Wall Drag Require Going Fully 3D?}
%==============================================================================

\begin{center}
\begin{longtable}{p{0.48\textwidth}p{0.48\textwidth}}
\toprule
\textbf{Scenario} & \textbf{Sufficient approach} \\
\midrule
Disruption-induced blanket flow, azimuthal symmetry, unresolved layers &
  2D axisymmetric + SM82 friction \\
\addlinespace
Same with Hartmann layers partially resolved at moderate Ha &
  2D axisymmetric with fine wall-normal mesh \\
\addlinespace
Duct flow with rectangular cross-section, conducting side walls &
  3D or reduced 2D duct model; axisymmetry is geometrically wrong \\
\addlinespace
Toroidal blanket asymmetry, field-ripple-driven 3D flow &
  Full 3D MHD \\
\addlinespace
Hartmann layer resolved to $\delta_H$ &
  3D with wall-normal anisotropic mesh \\
\addlinespace
Plasma--blanket feedback, disruption timescale &
  2D axisymmetric + TokaMaker coupling sufficient for a first study \\
\bottomrule
\end{longtable}
\end{center}

\textbf{Recommendation:} Prefer 2D axisymmetric with SM82 friction for all
first-generation blanket studies.  Full 3D becomes necessary only for
intrinsically 3D duct geometries, resolved Hartmann layers, or toroidal
asymmetry.

%==============================================================================
\section{Multi-Physics Coupling Infrastructure: Current State}
\label{sec:coupling_current}
%==============================================================================

\subsection{What Exists Today (Code-Confirmed)}

\subsubsection{TokaMaker $\to$ MUG: One-Way Equilibrium Export (Operational)}

\begin{itemize}
  \item \textbf{Fortran:} \texttt{gs\_save\_mug} in \texttt{src/physics/grad\_shaf.F90}
    (lines 4851--4939)
  \item \textbf{C wrapper:} \texttt{tokamaker\_save\_mug} in
    \texttt{src/python/wrappers/tokamaker\_f.F90} (line~1669)
  \item \textbf{Python:} \texttt{TokaMaker.save\_mug(filename)} in
    \texttt{TokaMaker/\_core.py} (lines~1905--1911) and
    \texttt{TokaMaker\_equil.save\_mug()} (lines~3012--3022)
  \item \textbf{Data transferred:} HDF5 file with mesh ($R$, $Z$), toroidal
    field ($B_R$, $B_T$, $B_Z$), pressure $P$, poloidal flux $\psi$
  \item \textbf{Mechanism:} Snapshot at a fixed time; one call per equilibrium state
\end{itemize}

\subsubsection{TokaMaker Live Field Interpolation (Operational)}

\begin{itemize}
  \item \textbf{Python class:} \texttt{TokaMaker\_field\_interpolator} in
    \texttt{TokaMaker/\_interface.py} (lines~253--291)
  \item \textbf{Core method:} \texttt{get\_field\_eval()} in \texttt{\_core.py}
    (lines~2606--2628)
  \item \textbf{Fortran bindings:} \texttt{tokamaker\_get\_field\_eval},
    \texttt{tokamaker\_apply\_field\_eval} (\texttt{\_interface.py}
    lines~161--167)
  \item \textbf{Fields available:} $\mathbf{B}$ (3-vector), $\psi$, $F$, $P$,
    $\nabla\psi$, $\nabla B_r$, $\nabla B_t$, $\nabla B_z$ at arbitrary $(R,Z)$
  \item \textbf{Status:} This is the right tool for live $\mathbf{B}_0(t)$
    injection into a low-Rm MUG solve, but no code currently connects it to MUG.
\end{itemize}

\subsubsection{ThinCurr Cross-Model Electromagnetic Coupling (Operational)}

\begin{itemize}
  \item \texttt{cross\_coupling()} in \texttt{ThinCurr/\_core.py}
    (lines~464--480): computes mutual inductance matrix $M$ between two
    ThinCurr models
  \item \texttt{cross\_eval()} (lines~482--498): computes flux on model~B due
    to currents on model~A
  \item \textbf{Tested:} \texttt{src/tests/physics/test\_ThinCurr.py} (line~267)
\end{itemize}

\subsubsection{ThinCurr Reduced Model (Operational)}

\begin{itemize}
  \item \texttt{build\_reduced\_model()} in \texttt{ThinCurr/\_core.py}
    (line~659): projects inductance/resistance matrices onto a reduced basis
  \item \texttt{ThinCurr\_reduced} class (lines~684--743): stores reduced
    operators for fast coupling
  \item Projection utilities in \texttt{ThinCurr/meshing.py}
    (lines~320--377): condense, expand, map between DOF representations
\end{itemize}

\subsubsection{TokaMaker--TORAX Temporal Coupling (Operational Template)}

\begin{itemize}
  \item \texttt{src/python/OpenFUSIONToolkit/TokaMaker/pulse\_design.py}
    implements an alternating-solve loop: TokaMaker equilibrium $\leftrightarrow$
    external transport code (TORAX)
  \item Coupling is temporal (profiles passed at each iteration); this is the
    architectural template for a TokaMaker--MUG coupling loop
\end{itemize}

\subsection{What Is Missing}

\begin{center}
\begin{longtable}{p{0.3\textwidth}p{0.65\textwidth}}
\toprule
\textbf{Missing piece} & \textbf{Description} \\
\midrule
MUG Python API & MUG is Fortran-only; no \texttt{ctypes}/\texttt{cffi}
  wrapper analogous to TokaMaker or ThinCurr.  This is the single largest
  structural gap --- without it, coupling must go through file I/O between
  separate processes. \\
\addlinespace
Live $\mathbf{B}_0$ injection into MUG & TokaMaker's
  \texttt{TokaMaker\_field\_interpolator} can supply $\mathbf{B}_0(R,Z,t)$
  at every MUG mesh node, but no code wires it to the MUG Fortran layer. \\
\addlinespace
$\mathbf{J}_{\mathrm{blanket}} \to$ TokaMaker & MUG's induced currents cannot
  feed back into the plasma equilibrium.  TokaMaker would need a
  ``blanket-current'' input analogous to its coil-current interface, implemented
  in \texttt{grad\_shaf.F90}. \\
\addlinespace
TokaMaker $\to$ ThinCurr spatial field map & TokaMaker's field interpolator
  provides point-wise evaluation, but no kernel maps TokaMaker's $(R,Z)$ fields
  onto ThinCurr's 3D triangular mesh for disruption wall-eddy-current analysis. \\
\addlinespace
ThinCurr $\to$ TokaMaker feedback & Wall eddy currents from disruptions cannot
  modify the plasma equilibrium; no coil-equivalent input from ThinCurr into
  TokaMaker exists. \\
\addlinespace
Multi-physics driver & No orchestration framework for simultaneous evolution of
  multiple physics domains beyond the point-coupling in \texttt{pulse\_design.py}. \\
\bottomrule
\end{longtable}
\end{center}

\subsection{Architecture of the Complete Coupling}

The full plasma--blanket--wall coupling described in the APS poster requires:

\begin{verbatim}
 ┌────────────────────────────────────────────────────────────────┐
 │                  Python Coupling Driver                        │
 │            (modelled after pulse_design.py)                    │
 └──────────┬────────────────────────────────┬───────────────────┘
            │  B0(R,Z,t)  [OPERATIONAL]      │  Jeq feedback
            ▼  save_mug() or interpolator    │  [MISSING]
 ┌──────────────────┐                        │
 │    TokaMaker     │ ◄──────────────────────┘
 │ (plasma equil.)  │
 └──────────────────┘
            │  B0(R,Z,t) via field interpolator
            ▼  [OPERATIONAL; NOT YET WIRED TO MUG]
 ┌──────────────────┐     J_blanket(R,Z,t)
 │   MUG xmhd_2d   │ ────────────────────►  TokaMaker
 │ (blanket flow,   │                        [MISSING]
 │  low-Rm mode)    │
 └──────────────────┘
            │  B0(t) on wall mesh
            ▼  [MISSING: TokaMaker→ThinCurr spatial map]
 ┌──────────────────┐
 │    ThinCurr      │
 │  (vessel wall    │ ──► TokaMaker coil equiv.  [MISSING]
 │   eddy currents) │
 └──────────────────┘
\end{verbatim}

The MUG Python API is the prerequisite for everything in the bottom half of
this diagram.

%==============================================================================
\section{Implementation Options for Wall Drag}
\label{sec:impl_options}
%==============================================================================

\subsection{Option A: Reduced Volumetric Hartmann Friction (\textbf{Recommended First PR})}

\subsubsection{Code Locations}

\textbf{Step 1 --- Add parameter to simulation object} \\
File: \texttt{src/physics/xmhd\_2d.F90}, inside \texttt{TYPE oft\_xmhd\_2d\_sim}
(lines 69--124), after line 92:

\begin{lstlisting}[caption={New field in oft\_xmhd\_2d\_sim (after line 92)}]
REAL(r8) :: drag_coeff = 0.d0  ! Hartmann/Shercliff friction coefficient
                                ! S_u = -drag_coeff * n * v
\end{lstlisting}

\textbf{Step 2 --- Add drag to nonlinear residual} \\
In \texttt{nlfun\_apply} (lines 717--1006), after momentum loop ($\sim$line 906):

\begin{lstlisting}[caption={Drag source term in nlfun\_apply after line~$\sim$906}]
!---Hartmann/wall friction source: S_u = -drag_coeff * n * v
IF (drag_coeff > 0.d0) THEN
  IF (cyl_flag) THEN
    DO k=1,3
      res_loc(jr,k+1) = res_loc(jr,k+1) &
        + self%dt * drag_coeff * basis_vals(jr) * n * vel(k) &
          * int_factor * coords(1)
    END DO
  ELSE
    DO k=1,3
      res_loc(jr,k+1) = res_loc(jr,k+1) &
        + self%dt * drag_coeff * basis_vals(jr) * n * vel(k) &
          * int_factor
    END DO
  END IF
END IF
\end{lstlisting}

\textbf{Step 3 --- Add drag to approximate Jacobian} \\
In \texttt{build\_approx\_jacobian} (lines 1010--1516), velocity self-coupling block
(lines 1232--1264):

\begin{lstlisting}[caption={Drag diagonal Jacobian in momentum block}]
IF (drag_coeff > 0.d0) THEN
  IF (cyl_flag) THEN
    jac_loc(k+1,k+1)%m(jr,jc) = jac_loc(k+1,k+1)%m(jr,jc) &
      + dt_fac * drag_coeff * n * basis_vals(jr) * basis_vals(jc) &
        * int_factor * coords(1)
  ELSE
    jac_loc(k+1,k+1)%m(jr,jc) = jac_loc(k+1,k+1)%m(jr,jc) &
      + dt_fac * drag_coeff * n * basis_vals(jr) * basis_vals(jc) &
        * int_factor
  END IF
END IF
\end{lstlisting}

\textbf{Step 4 --- Extract parameter} at top of both subroutines (after existing
material extractions at lines 748--756 and 1040--1049):
\begin{lstlisting}
drag_coeff = self%parent_sim%drag_coeff   ! nlfun_apply
drag_coeff = self%drag_coeff              ! build_approx_jacobian
\end{lstlisting}

\textbf{Step 5 --- Diagnostic:} add drag power
$P_\mathrm{drag} = \alpha \int n|\mathbf{v}|^2\,dV$ to history file.

\subsubsection{Input Convention}
\begin{verbatim}
&xmhd2d
  drag_coeff = 1.0e-3   ! dimensionless: alpha_H * tau_ref
/
\end{verbatim}

\subsection{Option B: Robin/Partial-Slip Velocity BC}

Boundary-integral traction term: $\oint \beta_s v_\parallel \phi \,dS$.
Requires new boundary-face quadrature loop --- more engineering effort, but
physically cleaner (wall-surface-only action).

\subsection{Option C: Region-Specific Liquid-Metal Closure}

Port \texttt{eta\_reg(:)} pattern from 3D module to support per-element
$\nu$, $\eta$, $\sigma$, $\rho$.  Enables multi-material domains.
Substantial effort; planned as a follow-on.

\subsection{Option D: TokaMaker / ThinCurr Coupling}

Discussed in Section~\ref{sec:coupling_roadmap}.  Presupposes Options A--C.

%==============================================================================
\section{Master Game Plan: Sequenced Pull Requests}
\label{sec:gameplan}
%==============================================================================

The following sequence moves from minimal first contribution to full
plasma--blanket coupling.  Each PR is independently useful and testable.

\begin{center}
\begin{longtable}{p{0.06\textwidth}p{0.25\textwidth}p{0.35\textwidth}p{0.24\textwidth}}
\toprule
\textbf{PR} & \textbf{Title} & \textbf{Scope} & \textbf{Delivers} \\
\midrule
1 & Scalar Hartmann drag & \texttt{drag\_coeff} field; residual + Jacobian in \texttt{xmhd\_2d.F90} ($\lesssim$80 lines); exponential decay test & First runnable LM drag \\
\addlinespace
2 & SM82 + buoyancy + floors & Ha-based $\alpha_H$, $\alpha_S$; \texttt{g\_accel}, Boussinesq; density/T floors; ohmic+viscous heating & PTOLEMY-comparable; magnetoconvection \\
\addlinespace
3 & Low-Rm $\varphi$ mode & New \texttt{use\_low\_rm} solver; Poisson solve for $\varphi$; remove Alfvénic constraint; insulating-wall Neumann BC; Hartmann channel test & HIMAG-comparable; $10^{2-3}\times \Delta t$ \\
\addlinespace
4 & Conducting-wall BC & Robin BC on $\psi$ (or $\varphi$); wall conductance ratio $c_w$; Hunt channel benchmark & Full conducting-wall physics \\
\addlinespace
5 & MUG Python API & \texttt{ctypes} wrapper for \texttt{xmhd\_2d}: \texttt{setup()}, \texttt{run\_n\_steps()}, \texttt{get\_fields()}, \texttt{set\_ic()} & Enables Python-level coupling \\
\addlinespace
6 & Live $\mathbf{B}_0$ injection & Wire TokaMaker \texttt{TokaMaker\_field\_interpolator} to MUG's low-Rm $\mathbf{B}_0$ input at each coupling step & Time-evolving background field \\
\addlinespace
7 & $\mathbf{J}_{\mathrm{blanket}} \to$ TokaMaker feedback & Add blanket-current input to \texttt{grad\_shaf.F90}; return induced flux to TokaMaker equilibrium & Bidirectional coupling \\
\addlinespace
8 & Full coupling driver & Python orchestration loop: TokaMaker $\to$ MUG $\to$ TokaMaker; convergence check; disruption example & APS poster capability \\
\bottomrule
\end{longtable}
\end{center}

\subsection{Timeline Guidance}

\begin{itemize}
  \item \textbf{PRs 1--2} can be done concurrently by a single developer in
    2--4 weeks.  No solver or architecture changes.
  \item \textbf{PR~3} (low-Rm mode) is the biggest single-PR effort
    ($\sim$400 lines of new Fortran + tests).  Estimate 3--6 weeks.
    This is the most physically impactful and removes the main computational
    bottleneck for liquid metals.
  \item \textbf{PR~4} (Robin BC) can follow PR~3 naturally; the $\varphi$
    formulation's Neumann BC is already correct for insulating walls, and
    Robin simply adds the $c_w$ term.  Estimate 1--2 weeks.
  \item \textbf{PR~5} (Python API) is primarily boilerplate following the
    TokaMaker/ThinCurr patterns.  Estimate 2--4 weeks.
  \item \textbf{PRs 6--8} build on each other quickly once PR~5 exists.
    Each is primarily Python orchestration.  Estimate 2--3 weeks each.
\end{itemize}

\subsection{Computational Cost Trajectory}

\begin{center}
\begin{tabular}{llll}
\toprule
\textbf{After PR} & \textbf{Bottleneck} & \textbf{Relative cost} & \textbf{HPC needed?} \\
\midrule
0 (current) & Alfvénic $\Delta t$, 7-field NL solve & Very high for LM & Possibly \\
1--2 & Same (drag adds nothing to cost) & Very high & Possibly \\
3 (low-Rm) & Flow timescale $\Delta t$, 5-field NL solve & Low--moderate & No \\
4 & Same + boundary integral & Low--moderate & No \\
5--8 (coupled) & Coupled iterations + TokaMaker equilibrium & Moderate & Modest HPC \\
\bottomrule
\end{tabular}
\end{center}

The low-Rm mode (PR~3) is the step that brings MUG into the same computational
class as PTOLEMY and HIMAG: desktop or modest cluster, not large HPC.

%==============================================================================
\section{Coupling Infrastructure Roadmap}
\label{sec:coupling_roadmap}
%==============================================================================

\subsection{Step 1: File-Based Coupling (Available Now)}

Use the existing \texttt{save\_mug()} to provide a TokaMaker equilibrium
snapshot to a standalone MUG Fortran simulation.  Iterate manually:
TokaMaker run $\to$ save HDF5 $\to$ MUG run $\to$ inspect output.
No new code needed; limited to static equilibria.

\subsection{Step 2: Live Field Injection via Python (After PR~5)}

With the MUG Python API in place, a Python script can:
\begin{enumerate}
  \item Call \texttt{TokaMaker.evolve(dt\_eq)} to advance the plasma equilibrium
  \item Call \texttt{tmaker\_interp.eval(mug\_mesh\_nodes)} to get
    $\mathbf{B}_0$ at all MUG nodes
  \item Inject $\mathbf{B}_0$ into MUG's low-Rm $\varphi$ solver
  \item Call \texttt{mug.run\_n\_steps(n)} to advance the blanket flow
  \item Repeat
\end{enumerate}
This is structurally identical to the TokaMaker--TORAX loop in
\texttt{pulse\_design.py}.

\subsection{Step 3: Bidirectional Feedback (After PR~7)}

Add the MUG-to-TokaMaker current feedback:
\begin{enumerate}
  \item MUG outputs $\mathbf{J}_{\mathrm{blanket}}(R,Z)$ (the induced current
    distribution in the liquid metal) via the Python API.
  \item TokaMaker ingests this as an axisymmetric current source in
    \texttt{grad\_shaf.F90}, treating it as an additional ``coil'' current.
  \item TokaMaker re-solves the equilibrium with the blanket current included.
  \item Iterate to convergence at each coupled timestep.
\end{enumerate}

\subsection{Step 4: ThinCurr Wall Integration (Optional Enhancement)}

Once the TokaMaker--MUG loop is running, add vessel-wall eddy currents:
\begin{enumerate}
  \item TokaMaker provides $\dot{\mathbf{B}}$ (field rate of change) during
    disruption.
  \item Interpolate onto ThinCurr's wall mesh using a new spatial projection
    kernel.
  \item ThinCurr evolves wall eddy currents.
  \item Wall currents feed back into TokaMaker as additional coil contributions.
\end{enumerate}
This completes the full disruption--wall--blanket--plasma coupling loop.

%==============================================================================
\section{Recommended First Pull Request}
%==============================================================================

\textbf{Title:} \texttt{xmhd\_2d: Add optional reduced Hartmann friction source term}

\textbf{Scope:}
\begin{enumerate}
  \item Add \texttt{drag\_coeff = 0.d0} to \texttt{oft\_xmhd\_2d\_sim} (line~$\approx$92)
  \item Extract at top of \texttt{nlfun\_apply} and \texttt{build\_approx\_jacobian}
  \item Add volumetric drag to momentum residual in \texttt{nlfun\_apply} ($\sim$12 lines)
  \item Add diagonal Jacobian to velocity block in \texttt{build\_approx\_jacobian} ($\sim$12 lines)
  \item Wire to Fortran namelist or XML input
  \item Add \texttt{test\_hartmann\_drag\_2d.F90/.py}: analytic exponential decay
  \item Add drag power diagnostic to history file
\end{enumerate}

\textbf{Do not include:} region-specific properties, Robin BC, TokaMaker coupling,
liquid-metal material constants.

\textbf{Estimated size:} $\lesssim$80 lines in \texttt{xmhd\_2d.F90}, $\sim$150 lines
for the new test.

%==============================================================================
\section{Testing and Validation Plan}
%==============================================================================

\subsection{Existing Baseline Tests (Must Still Pass)}

\begin{center}
\begin{longtable}{p{0.32\textwidth}p{0.22\textwidth}p{0.38\textwidth}}
\toprule
\textbf{Test} & \textbf{File} & \textbf{Pass criterion} \\
\midrule
2D Alfv\'{e}n wave & \texttt{test\_alfven\_2d.F90/py} & Mode frequency within tolerance; no change with \texttt{drag\_coeff=0} \\
2D sound wave & \texttt{test\_sound\_2d.F90/py} & Acoustic frequency within tolerance \\
Harris sheet & \texttt{harris\_sheet.F90} & Reconnection rate unchanged \\
\bottomrule
\end{longtable}
\end{center}

\subsection{New Unit Tests for Wall Drag}

\subsubsection{Test 1: Drag-Off Regression}
\texttt{drag\_coeff = 0}: bitwise-identical output to pre-PR baseline.

\subsubsection{Test 2: Analytic Exponential Decay}
Uniform $v_0$, $n = n_0$, $B = 0$, $\eta = \chi = D = \nu = 0$.
\begin{equation}
  \mathbf{v}(t) = \mathbf{v}_0 e^{-\alpha t}, \qquad
  \left|\frac{\|v(t_f)\|}{v_0} - e^{-\alpha t_f}\right| < 10^{-4}
\end{equation}

\subsubsection{Test 3: Kinetic Energy Monotone Decay}
$E_k(t_{i+1}) < E_k(t_i)$ at every stored timestep.

\subsubsection{Test 4: Sign Test}
$\mathbf{S}_u \cdot \mathbf{v} < 0$ at every quadrature point.

\subsubsection{Test 5: Parameter Parsing}
Negative \texttt{drag\_coeff} rejected or warned; zero value incurs zero overhead.

\subsubsection{Test 6: Jacobian Consistency}
Finite-difference Jacobian check with nonzero \texttt{drag\_coeff}.

\subsubsection{Test 7: Cylindrical Mode Decay}
Repeat Test~2 with \texttt{cyl\_flag = .TRUE.}; same analytic rate recovered.

\subsection{Post-PR Validation Benchmarks}

\begin{center}
\begin{longtable}{p{0.3\textwidth}p{0.65\textwidth}}
\toprule
\textbf{Benchmark} & \textbf{Description} \\
\midrule
Hartmann channel (analytic) & Quasi-2D pressure-flow relation with
  \texttt{drag\_coeff} = analytic $\alpha_H$; compare to Hunt (1965). \\
SM82 duct flow & Shercliff/Walker duct flow; pressure-drop comparison. \\
Damped Alfv\'{e}n + drag & Mode with both Alfv\'{e}nic oscillation and drag decay
  $\exp(-\alpha t/2)$; analytic damping formula. \\
Rayleigh--B\'{e}nard + $B$ & Magnetoconvection onset; compare to Burr \& Müller (2001). \\
Disruption blanket transient & Drag on/off; confirm peak velocity reduction. \\
\bottomrule
\end{longtable}
\end{center}

%==============================================================================
\section{Risks and Unknowns}
%==============================================================================

\begin{enumerate}[label=\textbf{R\arabic*.}]
  \item \textbf{Branch status of APS poster code.} May exist on a private branch.
    \textit{Action: ask Sophia Guizzo and Chris Hansen.}

  \item \textbf{Nondimensionalisation of \texttt{drag\_coeff}.} Normalisation
    tables are absent from code comments.
    \textit{Action: define explicitly relative to \texttt{den\_scale}, \texttt{m\_i},
    and reference time; enforce in tests.}

  \item \textbf{Density-weighted vs.\ mass-weighted drag.} Jacobian must match
    residual exactly.
    \textit{Action: finite-difference Jacobian check (Test~6).}

  \item \textbf{Region-specific drag.} Without regional material map, drag
    applies everywhere on the mesh.
    \textit{Action: document; plan as PR~E.}

  \item \textbf{Linear mode cylindrical restriction.} \texttt{run\_lin\_simulation}
    aborts for cylindrical (line 378--379).
    \textit{Action: add drag to linear residual in a follow-on PR.}

  \item \textbf{Low-Rm mode requires $\mathbf{B}_0$ as input.} When MUG does
    not evolve its own field, the background field must come from TokaMaker or
    an analytic prescription.  This is a design requirement, not a limitation.

  \item \textbf{Coupling convergence.} The alternating TokaMaker--MUG iteration
    may converge slowly for strongly coupled disruption events.
    \textit{Action: study convergence rate; consider sub-stepping MUG within each
    TokaMaker step.}

  \item \textbf{MUG Python API complexity.} TokaMaker's API wraps a complex
    Fortran object; the same for MUG may require careful memory management in
    the \texttt{ctypes} layer.
    \textit{Action: follow ThinCurr API as template; scope as a dedicated sprint.}
\end{enumerate}

%==============================================================================
\appendix
\section{Commands Run}
%==============================================================================

\begin{verbatim}
git status --short; git branch --show-current; git rev-parse HEAD
git cat-file -t 5faae9a; git show --stat 5faae9a
git show --name-status 5faae9a; git rev-parse 5faae9a^
git diff --stat 5faae9a^ 5faae9a; git diff --name-status 5faae9a^ 5faae9a

rg -n "liquid|lithium|lead|blanket|PbLi|TokaMaker|tokamaker|axisym|
  Hartmann|Shercliff|drag|friction|thin.wall|conducting.wall" \
  . --include="*.F90" --include="*.py" --include="*.in" \
  --include="*.xml" --include="*.md" -l

find . -iname "*mug*" -o -iname "*xmhd*" -o -iname "*blanket*" \
  -o -iname "*liquid*" -o -iname "*tokamaker*"

# Line-range reads in xmhd_2d.F90:
#   1-50, 69-135, 137-180, 860-960, 1010-1070, 1232-1270, 1530-1620

grep -n "def get\|def save\|def set\|def compute\|def solve" \
  src/python/OpenFUSIONToolkit/TokaMaker/_core.py

grep -n "cross_coupling\|cross_eval\|build_reduced\|ThinCurr_reduced" \
  src/python/OpenFUSIONToolkit/ThinCurr/_core.py

grep -n "coupling\|coupled\|transfer\|exchange" \
  src/python/OpenFUSIONToolkit/TokaMaker/pulse_design.py | head -30

find src/python/OpenFUSIONToolkit/ThinCurr -name "*.py" | sort
find src/examples/TokaMaker -type f | sort
\end{verbatim}

%==============================================================================
\section{Relevant Files, Modules, Subroutines, and Line Numbers}
%==============================================================================

\begin{center}
\begin{longtable}{p{0.52\textwidth}p{0.1\textwidth}p{0.32\textwidth}}
\toprule
\textbf{File / Symbol} & \textbf{Lines} & \textbf{Role} \\
\midrule
\texttt{src/physics/xmhd\_2d.F90} & 1--1812 & 2D MHD module (new in 5faae9a) \\
\quad \texttt{oft\_xmhd\_2d\_sim} type & 69--124 & Main simulation type \\
\quad \texttt{nu}, \texttt{eta}, \texttt{chi} fields & 83, 92, 82 & Material coefficients \\
\quad \texttt{cyl\_flag} & 71 & Cylindrical R-Z flag \\
\quad velocity BC arrays & 93--99 & Dirichlet BC pointers \\
\quad \texttt{run\_simulation} & 137--346 & NL time advance \\
\quad \texttt{run\_lin\_simulation} & 350--526 & Linear time advance \\
\quad \texttt{nlfun\_apply} & 717--1006 & Nonlinear residual \\
\quad viscosity in \texttt{nlfun\_apply} & 880--881, 903--904 & $\nu\nabla^2\mathbf{v}$ \\
\quad \textbf{drag injection point} & $\sim$906 & \textbf{After momentum DO k loop} \\
\quad \texttt{build\_approx\_jacobian} & 1010--1516 & 7$\times$7 block Jacobian \\
\quad \textbf{velocity self-coupling block} & 1232--1264 & \textbf{Jacobian drag injection point} \\
\quad \texttt{setup\_bc} & $\approx$1602--1611 & Default Dirichlet BCs \\
\midrule
\texttt{src/physics/xmhd.F90} & --- & 3D MHD module \\
\quad \texttt{nu\_par}, \texttt{nu\_perp} & 334--335 & Anisotropic viscosity \\
\quad \texttt{eta\_reg(:)} & 332 & Regional resistivity \\
\quad \texttt{eta\_temp}, \texttt{g\_accel} & 331, --- & T-dep.\ resistivity, gravity \\
\midrule
\texttt{src/physics/thin\_wall.F90} & --- & Thin-wall eddy current \\
\texttt{src/physics/thin\_wall\_solvers.F90} & --- & ThinCurr solvers \\
\texttt{src/physics/thin\_wall\_hodlr.F90} & --- & HODLR accelerator \\
\midrule
\texttt{src/physics/grad\_shaf.F90} & 4851--4939 & \texttt{gs\_save\_mug}: equilibrium export \\
\midrule
\texttt{src/python/.../TokaMaker/\_core.py} & 1905--1911 & \texttt{save\_mug()} Python call \\
 & 2606--2628 & \texttt{get\_field\_eval()} interpolator \\
 & 3012--3022 & \texttt{TokaMaker\_equil.save\_mug()} \\
\texttt{src/python/.../TokaMaker/\_interface.py} & 161--167 & Fortran field eval bindings \\
 & 253--291 & \texttt{TokaMaker\_field\_interpolator} \\
\texttt{src/python/.../TokaMaker/pulse\_design.py} & 230+ & TokaMaker--TORAX coupling loop \\
\midrule
\texttt{src/python/.../ThinCurr/\_core.py} & 464--480 & \texttt{cross\_coupling()} \\
 & 482--498 & \texttt{cross\_eval()} \\
 & 659 & \texttt{build\_reduced\_model()} \\
 & 684--743 & \texttt{ThinCurr\_reduced} class \\
\texttt{src/python/.../ThinCurr/meshing.py} & 320--377 & Projection utilities \\
\midrule
\texttt{src/python/wrappers/tokamaker\_f.F90} & 1669 & \texttt{tokamaker\_save\_mug} C wrapper \\
\midrule
\texttt{src/examples/MUG/harris\_sheet/harris\_sheet.F90} & 1--280 & 2D reconnection example \\
\texttt{src/examples/MUG/harris\_sheet/oft.in} & 23--49 & Input parameters \\
\texttt{src/tests/physics/test\_alfven\_2d.F90} & --- & Alfv\'{e}n wave test \\
\texttt{src/tests/physics/test\_sound\_2d.F90} & --- & Acoustic wave test \\
\texttt{src/tests/physics/test\_ThinCurr.py} & 267 & ThinCurr cross\_eval test \\
\bottomrule
\end{longtable}
\end{center}

%==============================================================================
\section{Unresolved Questions for Sophia / Maintainers}
%==============================================================================

\begin{enumerate}[label=\textbf{Q\arabic*.}]
  \item \textbf{Branch for APS poster code.}  Is the plasma--blanket coupling
    on a private/development branch?  What is blocking merge?

  \item \textbf{Normalisation convention.}  What is the reference time, density,
    length, and field normalisation for blanket simulations?
    Needed to correctly interpret \texttt{drag\_coeff} in dimensional units.

  \item \textbf{Liquid-metal material parameters.}  What values of $\nu$, $\eta$,
    $\chi$, $\sigma$ are used for PbLi or liquid lithium in the APS poster
    simulations, in \texttt{xmhd\_2d} normalised units?

  \item \textbf{Regional resistivity in 2D.}  Is there a plan to port
    \texttt{eta\_reg(:)} (3D \texttt{xmhd.F90}) to \texttt{xmhd\_2d.F90}?

  \item \textbf{Linear mode cylindrical restriction.}  Line 378--379 aborts
    for cylindrical in linear mode.  Planned limitation or temporary omission?

  \item \textbf{Coupling loop orchestration.}  Is there a Python driver (not in
    \texttt{main}) alternating TokaMaker and MUG solves?

  \item \textbf{Hartmann number estimates.}  Characteristic Ha for the blanket
    geometries of interest (ITER PbLi, liquid lithium)?

  \item \textbf{Low-Rm validity.}  Is Rm $\ll 1$ confirmed for the targeted
    blanket conditions, or are there regimes where induced-field evolution
    matters?

  \item \textbf{MUG Python API priority.}  Is there appetite among OFT
    maintainers to add a Python API for MUG, given that TokaMaker and ThinCurr
    already have one?  This is the prerequisite for all Python-level coupling.

  \item \textbf{Thin-wall coupling for 2D.}  Is there a 2D analogue of
    \texttt{thin\_wall.F90} planned, or should the conducting-wall effect be
    absorbed into an effective Ha via the wall conductance ratio $c_w$?
\end{enumerate}

\end{document}
